% Standalone solution dossier for ex:a5-neal.
\documentclass[../main.tex]{subfiles}
\begin{document}

% === SOLUTION HEADER (proof provenance) =====================================
% ledger-nodes: lem:a5-pi-exp-tail; ex:a5-neal
% refines: lem:a5-pi-exp-tail, ex:a5-neal
% bounded_by: obs:heavy-tail-no-classical
% evidence_run: none
% checked_by: agent
% author: /root/a4_a5
% proof_completion: /root/review_a4_a5
% dossier_editor: /root/review_a4_a5
% reviewer: /root/cross_review
% review: research/reviews/2026-08-21-a5-funnel-second-agent-audit.md
% date: 2026-08-21
% ============================================================================

\section*{Solution: Neal's Gaussian-log-scale funnel}

\begin{lemma}[Poincar\'e implies an exponential moment]
If a probability law satisfies a Poincar\'e inequality with finite constant $C$, every real
$1$-Lipschitz $f\in L^2$ has
$\E e^{c|f-\operatorname{med}f|}<\infty$ for every
$0<c<(\log2)/(2\sqrt C)$.
\end{lemma}

\begin{proof}
Let $m$ be a median, $a=2\sqrt C$, and, for $t\ge0$, set
$g=\min\{(f-m-t)_+,a\}$. If $p_t=\P(f>m+t)$, the independent-copy formula for variance gives
\[
 \Var(g)\ge a^2\P(f\ge m+t+a)\P(f\le m+t)
 \ge\tfrac12a^2p_{t+a}.
\]
Since $|\nabla g|\le\mathbf1_{\{m+t<f<m+t+a\}}$, Poincar\'e gives
$p_{t+a}\le(2C/a^2)p_t=\tfrac12p_t$. Iteration gives a geometric upper tail on intervals of
length $a$. Applying the same argument to $-f$ gives the lower tail. Integrating these two tails
proves the claimed exponential moment.
\end{proof}

\begin{proposition}
For $s>0$, let $u\sim N(0,s^2)$ and $z\sim N(0,1)$ be independent, and put $\theta=e^uz$.
In Euclidean
coordinates $(u,\theta)$ the Poincar\'e constant is infinite, while in coordinates $(u,z)$ it is
$\max\{s^2,1\}$. Neither coordinate map is globally Lipschitz.
\end{proposition}

\begin{proof}
The coordinate $\theta$ is one-Lipschitz on the Euclidean $(u,\theta)$ space and belongs to
$L^2$, since $\E\theta^2=e^{2s^2}$. Yet, for every $c>0$, independence and the event
$\{|z|\ge1,u>0\}$ give
\[
\E e^{c|\theta|}\ge\P(|z|\ge1)\E[e^{ce^u}\mathbf1_{\{u>0\}}]=\infty.
\]
The last Gaussian integral diverges because
$ce^u-u^2/(2s^2)\to+\infty$. If an exponential moment about the median $m$ were finite, then
$e^{c|\theta|}\le e^{c|m|}e^{c|\theta-m|}$ would make the uncentred moment finite. The lemma
therefore rules out a finite centered Poincar\'e constant. The non-centred law is a
product Gaussian, whose constant is $\max\{s^2,1\}$ by tensorization. The Jacobians of
$(u,z)\mapsto(u,e^uz)$ and $(u,\theta)\mapsto(u,e^{-u}\theta)$ are both unbounded.
\end{proof}

\paragraph{Audit state.}
The truncation lemma and funnel argument received a second independent agent audit after the
$s>0$ hypothesis and all integrability steps were made explicit.

\end{document}
